% Pass: Opus 5 (claude-opus-5), 11 September 2026 — first pass, unchecked
% against the views by a human.
%
% WHAT THESE TWENTY VIEWS HOLD. The head of the correspondence volume:
% BnF foliation 1 to 16, all of it letters, none of it mathematics set out as
% mathematics. One piece is Germain's own — a heavily corrected draft of a
% letter on the theory of elastic surfaces (ff. 1r-2r, views 5-6), unsigned,
% addressed « Monsieur », dated « ce 18 juillet » with no year, naming Savart
% and Ampère and asking its reader for his judgement on a « petit écrit » she
% has just had printed. Everything else is received: a bookseller's note
% (Bernard, quai des Augustins) about presenting « le citoyen Cousin » to her
% mother, ff. 3r-3v; two letters of Cauchy, f. 4r (1821, sending his own
% volume on algebraic analysis) and f. 6r (1826); two of Delambre as perpetual
% secretary, ff. 8r-8v (1816, on tickets for a public session) and f. 10r
% (1821, the Académie's thanks for « Recherches sur la théorie des surfaces
% élastiques »); and three of Fourier, f. 12r (1823, the reserved seat at the
% centre of the hall), f. 14r (1823, a private note, not on the Académie's
% paper), f. 16r (1824, « Mémoire sur l'emploi de l'épaisseur dans la théorie
% des surfaces élastiques » sent to Laplace, Prony and Poisson).
%
% HOW THE VIEWS ARE LAID OUT. Each view photographs an opening, not a leaf:
% the left half is the verso of the preceding folio, the right half the recto
% of the next, and the BnF's pencil numbers that right-hand recto. \folio{}
% therefore carries the right half's number; where the left half bears text,
% a \note{} says which verso it is. Views 1-4 are the microfilm target, the
% board, the marbled cover and the flyleaf with the old shelfmarks; views 9,
% 11, 13, 15, 17 and 19 are blank leaves showing only the ink of the other
% side. None of the seven appears below; the gap in the numbering is the
% record.
%
% THE HAND. The draft of ff. 1r-2r is fast, with the spelling of her
% generation (seroit, avoit, connoissance, tems, différens) and with a habit
% that a reader must not mistake: she finishes a great many words with a long
% horizontal flourish that looks exactly like a deletion stroke. Only
% cancellations that carry an interlinear replacement, or that strike a whole
% phrase twice, are read here as \struck{}. Where the rewriting is too tangled
% to give one line of text, the final reading stands in the body and the
% earlier one is recorded in a \note{}.
%
% NOTATION. The ordinal abbreviations are set on the line as the nearest
% standard forms — « 1° », « n° », « 3me », « 1er », « Bon » for the baron's
% title — where the leaves raise the ending above the line.
%
% The section headings are the transcriber's; not one of them is written on
% the leaves.
\documentclass[11pt,a4paper]{article}
% The preamble shared by every transcription of the mathematicians' manuscripts in Gallica.
%
% It rests on one idea: **uncertainty compounds, it does not resolve.** A
% transcription that smooths over a doubtful word has destroyed the only thing
% it was made for — a reader must be able to tell, at every point, what was on
% the page and what was supplied. The apparatus macros below are therefore the
% critical apparatus, not decoration.
%
% The same commands are recognised by `scripts/render.mjs`, which produces the
% site's reading view, and by `scripts/tei.mjs`. A macro added here must be
% added there too, or rendering fails loudly — which is the wanted behaviour.
%
% Comments here are English, like the rest of the repository. The documents
% this preamble sets are French, and that is deliberate: see the skills.

\ProvidesPackage{germain}[2026/09/15 Transcriptions of mathematicians' manuscripts in Gallica]

\usepackage{amsmath,amssymb,amsthm}
\usepackage{mathrsfs}
% Kept from the parent preamble so the renderer's diagram support stays
% available; these manuscripts draw no commutative diagrams, and nothing here uses it.
\usepackage{tikz-cd}
\usepackage[english,french]{babel}
\usepackage{fontspec}

% --- Greek ------------------------------------------------------------------
% The text face stays fontspec's default, Latin Modern, which has no Greek:
% XeTeX drops every glyph it lacks with a line in the log and nothing on the
% page, and the Greek words the manuscripts quote vanished from the PDFs. So
% the Greek and Greek Extended blocks get a XeTeX character class of their own,
% and entering or leaving a run of it switches the family to CMU Serif — the
% same Computer Modern design, with polytonic Greek — and back. Only the
% family changes: a Greek word in \textit or \textbf keeps its shape and series.
% The transitions to and from every other class carry the tokens that class
% has towards an ordinary letter, so French spacing before « ; » or inside
% « » still holds after a Greek word. No grouping, so no brace or \endgroup
% can come out unbalanced; maths is left alone, since XeTeX inserts nothing
% there. Kept identical in `exercices.sty`.
\makeatletter
\newfontfamily\germainGreekfont{cmun}[
  Extension=.otf, NFSSFamily=germaingreek,
  UprightFont=*rm, ItalicFont=*ti, SlantedFont=*sl,
  BoldFont=*bx, BoldItalicFont=*bi, BoldSlantedFont=*bl]
\newXeTeXintercharclass\germain@greek
\def\germain@greekrange#1#2{%
  \@tempcnta=#1\relax
  \loop
    \XeTeXcharclass\@tempcnta=\germain@greek
    \ifnum\@tempcnta<#2\advance\@tempcnta\@ne\repeat}
\germain@greekrange{"0370}{"03FF}
\germain@greekrange{"1F00}{"1FFF}
\def\germain@greekprev{\rmdefault}
\def\germain@greekin{%
  \edef\germain@greekprev{\f@family}\fontfamily{germaingreek}\selectfont}
\def\germain@greekout{\fontfamily{\germain@greekprev}\selectfont}
\def\germain@greekjoin{%
  \XeTeXinterchartokenstate=\@ne
  \@tempcnta=\z@
  \loop
    \ifnum\@tempcnta=\germain@greek\else
      \XeTeXinterchartoks\germain@greek\@tempcnta=\expandafter{%
        \expandafter\germain@greekout\the\XeTeXinterchartoks\z@\@tempcnta}%
      \XeTeXinterchartoks\@tempcnta\germain@greek=\expandafter{%
        \the\XeTeXinterchartoks\@tempcnta\z@\germain@greekin}%
    \fi
    \ifnum\@tempcnta<4095\advance\@tempcnta\@ne\repeat}
% babel-french sets its own classes, and saves and restores the interchar
% state, each time a language is selected: join again after it does.
\addto\extrasfrench{\germain@greekjoin}
\addto\extrasenglish{\germain@greekjoin}
\AtBeginDocument{\germain@greekjoin}
\makeatother

\usepackage{geometry}
\geometry{a4paper,margin=25mm,marginparwidth=28mm}
\usepackage{marginnote}
\usepackage{xcolor}
\usepackage[normalem]{ulem}
\setlength{\emergencystretch}{3em}
\usepackage{hyperref}
\hypersetup{colorlinks=true,linkcolor=black,urlcolor=black,pdfborder={0 0 0}}

% --- Batch metadata ---------------------------------------------------------
% Taken as they stand by the rendering script, which shows them at the head of
% the reading view. They fix what the file claims to cover. The macro names are
% the parent project's, so that the scripts stay shared; \folder here means the
% volume's slug — `fr-9115` — and \pages counts Gallica views.

\newcommand{\folder}[1]{\gdef\grothFolder{#1}}
\newcommand{\batch}[1]{\gdef\grothBatch{#1}}
\newcommand{\pages}[2]{\gdef\grothFirst{#1}\gdef\grothLast{#2}}
\newcommand{\foldertitle}[1]{\gdef\grothFoldertitle{#1}}

% The holder's own shelfmark, printed as they write it: « Français 9115 ».
\newcommand{\shelfmark}[1]{\gdef\grothShelfmark{#1}}

% The dating, copied verbatim from the holder's catalogue — never inferred
% here. « XIXe siècle » is the BnF's; a pass that dates a leaf from its content
% says so in a \note{}, as its own inference.
\newcommand{\dating}[1]{\gdef\grothDating{#1}}

% The status notice, printed in the title block and repeated as a diagonal
% watermark on every page of the PDF. The manuscripts are in the public
% domain; what this line declares is not a right but a quality — an
% unverified first pass by a machine. The skills write the line; a file
% without it gets no watermark, so leaving it out is visible in the output.
\newcommand{\watermark}[1]{\gdef\grothWatermark{#1}}

\gdef\grothFolder{?}\gdef\grothBatch{}\gdef\grothFirst{?}\gdef\grothLast{?}%
\gdef\grothFoldertitle{}\gdef\grothShelfmark{}\gdef\grothDating{}\gdef\grothWatermark{}

% --- The résumé -------------------------------------------------------------
\newenvironment{resume}
  {\small\setlength{\parskip}{0.45em}}
  {\par\medskip\hrule\medskip}

% --- Keywords ---------------------------------------------------------------
% In English, closing the modernised reading's résumé: the modern vocabulary
% under which to search for what the leaves construct. The manifest extracts
% them to tag the volume — this line is the single source of the tags.
\newcommand{\keywords}[1]{\par\smallskip\noindent
  {\footnotesize\textsc{Keywords} --- \textit{#1}}\par}

% --- The critical apparatus -------------------------------------------------

% The Gallica view number, in the margin. This is the anchor that lets a
% disagreement over a formula be settled by looking at *one* image rather than
% a volume of 750. The site uses it to turn the facsimile as the transcription
% is scrolled. A view is one scanned image: usually one side of a leaf.
\newcommand{\page}[1]{%
  \marginnote{\footnotesize\color{gray!70}\textsf{v.\,#1}}%
  \label{page:#1}\ignorespaces}

% The BnF's pencilled foliation, where the leaf carries it and a pass can read
% it: « 198r ». This is what the literature cites — « ff. 198r–208v » — and
% the one line that turns a citation into a view. Recorded, never inferred:
% a folio a pass cannot read is not written.
\newcommand{\folio}[1]{%
  \marginnote{\footnotesize\color{gray!50}\textsf{f.\,#1}}\ignorespaces}

% The modernised reading's coarser counterpart to \page{}: which views a
% section is drawn from.
\newcommand{\pagerange}[2]{%
  \marginnote{\footnotesize\color{gray!70}\textsf{v.\,#1--#2}}%
  \ignorespaces}

% Illegible: never guessed.
\newcommand{\ill}{\textcolor{red!60!black}{[\,\ldots\,]}}

% A doubtful reading: the word is given, and flagged as uncertain.
\newcommand{\uncertain}[1]{\uline{#1}}

% An editorial addition — a missing letter, an implied word.
\newcommand{\add}[1]{\textcolor{blue!50!black}{[#1]}}

% Struck out by the writer. What was crossed out is part of the document.
\newcommand{\struck}[1]{\sout{#1}}

% The transcriber's note, on the state of the page or the mathematics.
\newcommand{\note}[1]{\par\smallskip{\small\color{gray!60!black}\textit{[#1]}}\par\smallskip}

% A marginal note of hers, distinct from the body of the text.
\newcommand{\marginal}[1]{\par\smallskip\noindent\hspace{1em}%
  {\small\color{gray!60!black}#1}\par\smallskip}

% --- Title ------------------------------------------------------------------
% The title block is French, like the documents themselves.

\AddToHook{shipout/background}{%
  \ifx\grothWatermark\empty\else
    \begin{tikzpicture}[remember picture, overlay]
      \node[rotate=52, scale=3.4, align=center, text=gray!18,
            font=\sffamily\bfseries]
        at (current page.center) {\grothWatermark};
    \end{tikzpicture}%
  \fi}

\AtBeginDocument{%
  \begin{center}
    {\large\bfseries Manuscrits de mathématiciens (Gallica) ---
      \ifx\grothShelfmark\empty\grothFolder\else\grothShelfmark\fi}\\[2pt]
    {\itshape\grothFoldertitle}\\[4pt]
    {\small vues \grothFirst--\grothLast
      \ifx\grothBatch\empty\else\ (lot \grothBatch)\fi}\\[2pt]
    \ifx\grothDating\empty\else
      {\small\color{gray!60!black}Datation du catalogue : \grothDating}\\[2pt]
    \fi
    \ifx\grothWatermark\empty\else
      {\small\bfseries\color{red!45!gray}\grothWatermark}\\[2pt]
    \fi
    {\footnotesize\color{gray!60!black}Transcription établie d'après les images IIIF de Gallica.
     Source gallica.bnf.fr / Bibliothèque nationale de France.\\
     Les lectures incertaines sont soulignées, les illisibles notées [\,\ldots\,],
     les ajouts entre crochets ; \textsf{v.} numérote les vues de Gallica,
     \textsf{f.} la foliotation au crayon de la BnF.}
  \end{center}
  \vspace{1em}}

\endinput


\folder{fr-9118}
\shelfmark{Français 9118}
\batch{1}
\pages{1}{20}
\dating{1801-1900}
\watermark{Lecture automatique\\non vérifiée}
\foldertitle{Correspondance de Sophie GERMAIN avec les mathématiciens et savants Cauchy, Delambre, Fourier, Gauss, Le Gendre, d'Ansse de Villoison, etc.}

\begin{document}

\section{Brouillon d'une lettre sur la théorie des surfaces élastiques}

\page{5}\folio{1r}
\note{feuillet de la main qui corrige, non signé ; l'attribution à Sophie
Germain est l'inférence de cette lecture, tirée du contenu — l'auteur y parle
du « premier mémoire que j'ai publié » sur les courbures moyennes — et non
d'une signature. Le destinataire n'est pas nommé.}

ce 18 juillet

Monsieur

J'ai vu avec plaisir que mes nouvelles remarques ont été renvoyées à votre
jugement : aucun autre ne m'auroit paru aussi sûr.

J'ai suivi votre avis en imprimant je m'estimerai heureuse si vous prenez la
peine de me communiquer vos observations. Il y a dans ce petit écrit trois
choses qui me sembloient de quelqu'importance 1° la définition de la question
d'où résultent la connoissance des conditions de ce cas particulier du
mouvement des corps doués d'élasticité. On voit, sans que j'aie eu besoin de
le dire, que les nombreuses expériences de Monsieur Savart sont étrangères à
cette question, elles ne pourroient être expliquées qu'à l'aide d'une théorie
plus étendue et telle que celle dont vous vous êtes occupé.
\note{la première rédaction, biffée et corrigée en interligne, portoit « et
telle que celle sur laquelle vous avez déjà porté votre attention » ; « dont »
est écrit au-dessus de « sur laquelle », « vous êtes » au-dessus de « avez
déjà », « occupé » au-dessus de « porté votre attention »}
En effet si dans les expériences dont je parle on parvenoit à séparer les
différentes couches dont on peut concevoir que l'épaisseur soit composée,
chacune d'elles présenteroit des figures

\page{6}\folio{2r}
\note{la moitié gauche de la vue est le verso du feuillet 1, où le brouillon
se poursuit ; la moitié droite est le feuillet 2 recto}

particulières, l'épaisseur variroit aussi à raison du mouvement et en effet
l'expérience rend alors sensible les différences d'épaisseurs produites par le
mouvement.

J'avois prié Monsieur Ampère de vous demander où je pourrois trouver ce que
vous avez publié sur le cas général du mouvement des corps élastiques, il ne
l'a pas fait et je n'ai pu retrouver qu'un premier aperçu insuffisant à mon
instruction.

Autant que je puis me rappeler, ce que vous avez pris la peine de m'expliquer
vous-même, le mouvement des surfaces présenteroit le cas particulier où la
résultante des forces qui agiroient sur chacune des molécules seroit
perpendiculaire aux différens plans tangens.
\marginal{agissant \uncertain{suivant} toutes les directions possibles}
\note{la phrase est reprise en interligne d'une écriture beaucoup plus fine et
en partie illisible : « \ill{} être considéré comme composé et produit par des
forces qui \ill{} » ; la leçon donnée ci-dessus est celle de la ligne de base}
Je serois bien enchantée que vous voulussiez reprendre ce genre de recherches
et mon foible travail prendroit \uncertain{à mes yeux} une importance réelle
s'il pouvoit contribuer à y ramener votre attention.

Une seconde considération sur laquelle je voudrois avoir votre avis est celle
des courbures moyennes j'en avois déjà parlé dans le premier mémoire que j'ai
publié. Il arrive par rapport à la courbure ce qu'on observe dans une foule
d'autres manières d'être des corps. je veux dire que l'état réel des points
dont la position est également éloignée de ceux \struck{\ill{}}
\uncertain{auxquels} ils appartiennent, les manières d'être extrêmes donnent
l'état moyen du système. C'est ainsi que dans \uncertain{l'enceinte} la
température des points qui sont également éloignés de ceux qui
\uncertain{possèdent} le maximum et le minimum de température
\struck{ont la} \struck{n'est autre chose} est égale à la température moyenne
de la pièce. Ici il faudroit une égale quantité de chaleur pour faire changer
une température moyenne donnée, à l'égard des surfaces il faudroit des forces
égales pour faire changer une courbure moyenne donnée, en sorte que la
courbure de la sphère est toujours comparable à celle d'une surface de figure
quelconque. Il m'a paru que cette remarque pouvoit être de quelqu'utilité et
elle m'a servi à donner une forme assez simple à l'équation générale des
surfaces. Enfin cette équation même me semble incontestable et je mettrois
beaucoup d'intérêt à savoir ce que vous en pensez.

Peut-être, Monsieur, trouverez-vous que c'est abuser de votre complaisance
que \struck{d'avoir demandé} d'ajouter l'ennui de ce commentaire à celui de la
lecture de mon petit mémoire ; je n'aurois d'autre excuse que l'importance que
je mets à votre jugement, je vous prie Monsieur d'en agréer l'assurance en
même tems que celle de mon respect.
\note{la lettre s'arrête ici, sans signature ; le brouillon ne porte pas
d'année}

\section{Bernard, libraire, à la mère de Sophie Germain}

\page{7}\folio{3r}

Paris 4 novembre

Madame,

Le citoyen Cousin sollicite l'honneur de vous être présenté ainsi qu'à
Mademoiselle votre fille, si vous daignez l'agréer. j'attends vos ordres. il
espère que ses occupations lui permettront d'aller vous présenter ses respects
le huit novembre sur les six heures du soir, je vous prie de me faire dire si
cette heure ne vous est pas importune. nous aurions l'honneur de nous rendre
au moment qui vous seroit plus convenable. m. cousin se félicite d'avoir une
occasion de vous offrir, ainsi qu'à Mademoiselle Germain, l'hommage de son
respect et de

\page{8}\folio{4r}
\note{la moitié gauche de la vue est le feuillet 3 verso, où la lettre du
libraire s'achève ; la moitié droite est le feuillet 4 recto, une lettre de
Cauchy}

lui offrir toutes les facilités qui dépendront de lui dans la carrière des
sciences qu'elle cultive avec tant de succès. je serai bien flatté d'avoir pu
concourir à vous offrir quelque témoignage de mon zèle, et de ma sincère
admiration pour votre idole.

Recevez, Madame, je vous prie pour elle et pour vous l'assurance de mon
parfait attachement.

Bernard libraire quai des augustins n° 27.
\note{« le citoyen Cousin » place la lettre sous le Directoire ou le
Consulat ; l'année n'est pas écrite, et cette datation est l'inférence de
cette lecture, non une mention du feuillet}

\section{Augustin Cauchy, 1821}

\note{feuillet 4 recto, moitié droite de la vue 8}

Mademoiselle,

J'ai reçu l'ouvrage que vous avez eu la bonté de m'adresser, ouvrage que le
nom de son auteur et l'importance du sujet recommandent également à
l'attention des géomètres. Je n'ai pour le moment à vous offrir en revanche
qu'un volume dans lequel j'ai cherché à éclaircir les principales difficultés
de l'analyse algébrique. Veuillez bien l'agréer, je vous prie, avec l'hommage
de ma considération distinguée et de mes très humbles respects.

Augustin Cauchy

Paris ce \uncertain{24} juillet 1821
\note{le quantième est repris à la plume sur un premier chiffre et reste
douteux ; l'année est nette}

\section{Augustin Cauchy, 1826}

\page{10}\folio{6r}
\note{la moitié gauche de la vue est le feuillet 5 verso, qui porte
l'adresse : « À Mademoiselle / Mademoiselle Sophie Germain / Rue \ill{}
Braque, n° / À Paris »}

\uncertain{Sceaux Penthièvre} ce 23 juillet 1826

Mademoiselle,

J'ai reçu la lettre que vous m'avez fait l'honneur de m'écrire avec le mémoire
qui l'accompagnoit. Je vous remercie d'avoir bien voulu m'adresser un
exemplaire de cet ouvrage, que je relirai avec tout le soin que réclament et
l'importance du sujet et le mérite de l'auteur.

Agréez, je vous prie, l'hommage du respect avec lequel je suis

Mademoiselle

votre très humble et très obéissant serviteur

A. L. Cauchy

\section{Delambre, secrétaire perpétuel, 1816}

\page{12}\folio{8r}
\note{la moitié gauche de la vue est le feuillet 7 verso, qui porte l'adresse
et un cachet de poste : « À Mademoiselle / Mademoiselle Sophie Germain / rue
de Braque, N° 4 / au marais »}

Institut de France. Classe des Sciences Physiques et Mathématiques.

Paris, le 6 Janvier 1816

Le Secrétaire perpétuel pour les Sciences

M. Delambre a l'honneur de présenter ses hommages à Mademoiselle Germain et de
lui envoyer deux billets d'Institut présumant bien que ses amis lui en
demanderont plus qu'elle n'en aura à distribuer. Si comme il le suppose elle
en a reçu hier ou aujourd'hui ; mais M. Delambre ayant appris par
\uncertain{Mde de Sédillot} que Mademoiselle Germain n'en avoit pas encore
reçu avant hier soir, il craint qu'il n'y ait eu quelqu'oubli, et la prie,
dans ce cas, d'avoir recours à lui, parceque, les billets imprimés étant
épuisés, il peut y suppléer par un billet à la main pour autant de personnes
qu'il conviendra à Mademoiselle Germain de lui en indiquer ; M. Delambre
désireroit bien qu'elle se rendît elle-même à la séance, il

\page{13}\folio{9r}
\note{la moitié gauche de la vue est le feuillet 8 verso, où la lettre
s'achève ; la moitié droite, feuillet 9 recto, est blanche}

auroit le plus grand plaisir à lui faire son compliment et lui renouveller
l'assurance de sa respectueuse considération.

Rue du petit Bourbon St Sulpice n° 7 maison du notaire

\section{Delambre, secrétaire perpétuel, 1821}

\page{14}\folio{10r}
\note{la moitié gauche de la vue est le feuillet 9 verso, qui porte
l'adresse : « À Mademoiselle / Mademoiselle Sophie Germain / Rue Ste Croix de
la Bretonnerie n° / la 3me porte cochère à droite après la rue \ill{}, avant
la rue \ill{} / À Paris »}

Institut de France. Académie Royale des Sciences.

Paris, le 23 juillet 1821.

Le Secrétaire perpétuel de l'Académie. À Mademoiselle Sophie Germain

Mademoiselle

L'Académie a reçu avec le plus grand intérêt l'ouvrage que vous avez bien
voulu lui adresser et qui est intitulé : \emph{Recherches sur la théorie des
surfaces élastiques} que vous venez de publier. Elle me charge de vous
remercier, en son nom, de l'envoi de ce mémoire intéressant qu'elle a fait
déposer honorablement dans la bibliothèque de l'Institut, et de vous exprimer
sa reconnaissance de cette nouvelle preuve que vous lui donnez de vos talents.

Agréez je vous prie, Mademoiselle, l'hommage de mon respect.

Delambre

\section{Fourier, secrétaire perpétuel, 1823}

\page{16}\folio{12r}
\note{la moitié gauche de la vue est le feuillet 11 verso, qui porte
l'adresse : « À Mademoiselle / Mademoiselle Sophie Germain / Rue de Braque
N° 4 / \ill{} / À Paris »}

Institut de France. Académie Royale des Sciences.

Paris, le 30 Mai 1823.

Le Secrétaire perpétuel de l'Académie.

Mademoiselle,

J'ai l'honneur de vous prévenir que toutes les fois que vous vous proposerez
d'assister aux séances publiques de l'Institut vous y serez admise dans l'une
des places réservées au centre de la salle. L'académie des sciences désire
témoigner par cette distinction tout l'intérêt que lui inspirent vos ouvrages
mathématiques et spécialement les savantes recherches qu'elle a couronnées en
vous décernant un de ses grands prix annuels.

Agréez, Mademoiselle, l'hommage de mon respect.

fourier.

Mademoiselle Sophie Germain.

\section{Fourier, billet particulier, 1823}

\page{18}\folio{14r}
\note{feuillet 14 recto ; papier ordinaire, sans l'en-tête gravé de
l'Académie, et d'une écriture beaucoup plus rapide que celle des lettres
officielles}

1er juin 1823.

J'ai l'honneur de me rappeler au souvenir et à la bienveillance de
mademoiselle germain. je désirois depuis longtems me présenter chez elle, mais
des occupations urgentes m'en ont détourné. je lui envoye cy-joint 1° une
lettre officielle 2° un billet du centre pour la personne qui l'accompagnera.
si mademoiselle germain ne se propose pas d'assister à la séance je la prie de
disposer du billet comme elle le jugera convenable et s'il étoit nécessaire
j'en pourrois remettre un \struck{ou plusieurs autres} mais non du centre.

hélas je \uncertain{devrois} bien plutôt garder tous les billets. je suis
condamné à causer au public un grand ennui et je vais paroître demain comme
une foible lueur au milieu d'un feu d'artifice. mais je suis résigné à toutes
les \uncertain{comparaisons} possibles. il m'a paru convenable de prendre dès
le début un ton grave et simple que je puis conserver et de m'abstenir de
toute prétention à des succès que je ne pouvois pas obtenir et que je ne
désire point. ce que je désire surtout c'est de conserver l'estime et le
souvenir de mademoiselle germain. je la prie de recevoir l'expression de mon
respect.

fourier.

\section{Fourier, secrétaire perpétuel, 1824}

\page{20}\folio{16r}

Institut de France. Académie Royale des Sciences.

Paris, le 15 Mars 1824.

Le Secrétaire perpétuel de l'Académie. À Mademoiselle Sophie Germain.

Mademoiselle,

L'Académie a reçu l'ouvrage que vous avez bien voulu lui adresser et qui est
intitulé : \emph{Mémoire sur l'emploi de l'épaisseur dans la théorie des
surfaces élastiques}. J'ai l'honneur de vous prévenir que ce mémoire a été lu
dans la séance du lundi 9 mars, et qu'il a été renvoyé à l'examen d'une
Commission composée de MM. Delaplace, De Prony et Poisson. Je m'empresserai de
mettre à votre disposition le rapport que l'académie aura adopté.

Agréez, Mademoiselle, l'expression de mon respect.

Bon fourier.
\note{« Delaplace » est écrit ainsi sur le feuillet ; la signature porte le
titre de baron abrégé}

\end{document}
